\section{海洋を水・運動量・物質の輸送系として理解する}
\label{theory:ocean:overview}
海洋は塩水の巨大な貯蔵庫であると同時に、風と密度差により運動し、熱と炭素と栄養塩を運ぶ系である。海水の組成、循環、生物による物質移動、海底での堆積を一つの収支としてつなぐ。津波では同じ海水でも長波の動的応答が中心となる。

\subsection{水の割合は何を分母にするか}
\label{theory:ocean:reservoir}
海水の体積は地球に存在する水の総体積の代表的推計で約96.5\%を占める。地球表面のうち海が占める面積の割合とは異なる。水量を時間不変の厳密値とせず、氷期の海水準変動なども含む貯蔵量と理解する。\cite{usgs_ocean_water_volume}

\subsection{35‰と塩類組成の85\%}
\label{theory:ocean:salinity}
原典の外洋の平均塩分は35‰である。質量比として\(35\,\mathrm g/1000\,\mathrm g=0.035\)なので約3.5\%となり、35\%ではない。海水全体に占める塩類と、その塩類中に占めるCl\(^{-}\)・Na\(^{+}\)の約85\%は異なる分母である。塩分の地域差とこの成分比を混同しない。\cite{noaa_sea_salt}
保存的なイオンの質量\(m_i,m_j\)を変えず水だけを蒸発又は加える理想モデルでは、濃度\(c_i=m_i/M\)の分母\(M\)が共通に変わり、比\(c_i/c_j=m_i/m_j\)は保たれる。これは保存性を前提とした収支である。反応、河川由来の異なる組成、熱水、微量の栄養塩には同じ比を保証しない。実用塩分という観測尺度と質量による絶対塩分の厳密な差まで、原典の概数に読み込まない。

\subsection{表層の風成循環と地衡流}
\label{theory:ocean:gyres}
風応力は海へ運動量を渡し、コリオリ力と圧力勾配が大規模な流れの方向に関わる。密度一定、水平流速\((u,v)\)、鉛直上向き、コリオリ係数\(f=2\Omega\sin\phi\)とする地衡流近似は
\[
 -fv=-\rho^{-1}\partial_xp,\qquad
 fu=-\rho^{-1}\partial_yp
\]
である。海面が高い領域の外向き圧力勾配に釣り合う循環は、\(f>0\)の北半球では時計回り、\(f<0\)の南半球では反時計回りとなる。これを亜熱帯環流の典型像へ対応させる。亜寒帯環流の海面の高低と回転方向は別であり、南半球の全ての渦を同じ向きとしない。\cite{noaa_subtropical_ecco}
赤道付近で\(f\)が小さいとこの単純な地衡流式だけでは応答を表せない。表層流が単に風向きと一致するわけでもない。

\subsection{密度差と深層循環の時間尺度}
\label{theory:ocean:overturning}
冷却と塩分増加は一般に海水の密度を増し、高密度水の沈降に関わる。北大西洋深層水や南極底層水など複数の水塊の形成、混合、湧昇を通じて海洋の深部と表層はつながる。千年程度という模式的な循環時間を、一つの流体粒子の一意の閉軌道や全水塊の絶対的最大年齢としない。「最大2000年」という原典は定義の追加確認が必要である。\cite{noaa_overturning}

\subsection{栄養塩：表層で消費され深部で再生される}
\label{theory:ocean:nutrients}
光の届く表層では植物プランクトンがリンやケイ素などを取り込む。沈降する有機物や殻が分解・溶解して水へ戻すと、深部の栄養塩が増える。典型的な表層低・深層高という向きは、この生物による移動の収支から理解できる。湧昇、季節変動、水塊の年齢、種類ごとの溶解を含む実分布は単調ではない。\cite{noaa_nutrient_observations}
濃度を\(C\)、流速を\(\boldsymbol u\)、拡散係数を\(D\)、再生と消費を\(R,U\)とする概念モデルは
\[
 \partial_t C+\boldsymbol u\cdot\nabla C
 =\nabla\cdot(D\nabla C)+R-U
\]
である。深さの関数だけでなく、移流と生物作用を含む点が保存的な主要塩類との違いである。

\subsection{炭酸塩の飽和・溶解速度・補償深度}
\label{theory:ocean:carbonate}
炭酸塩の飽和度は、所定の温度・圧力・塩分での溶解度積に対するイオン積の比で表せる。熱力学的に溶けやすいことと、一定時間にどれだけ溶けるかは別である。実際の速度は飽和度、粒子、境界層と海域に関わり、「深海ほど必ず速い」とする単調則は採らない。\cite{feely_carbonate2004}
炭酸塩補償深度は、海底への炭酸塩の供給と溶解の釣合いで残存が変わる深度である。したがって「深海」という名称だけでは石灰質軟泥の有無を確定できず、補償深度より上の深海底では残り得る。飽和深度、溶解が顕著になる深度、補償深度を同一の面としない。\cite{noaa_ccd_definition,feely_carbonate2004}

\subsection{団塊とクラストを区別する}
\label{theory:ocean:nodules}
マンガン団塊は深海底の核の周囲にゆっくり成長する鉄・マンガン酸化物の集合である。海山表面のマンガンクラストとは形態と位置が違う。成長が遅いため、速い堆積で埋まらず海底面に保たれる環境が形成・残存に関わるというモデルで理解する。堆積の全域に一律の速度を与えない。\cite{jamstec_deep_resources}

\subsection{CO2の吸収と海水の酸性化}
\label{theory:ocean:carbon}
溶けたCO2は炭酸水素イオンなどへ分配される。簡略化した反応は
\[
 \mathrm{CO_2(aq)+H_2O\rightleftharpoons H^++HCO_3^-},\qquad
 \mathrm{HCO_3^-\rightleftharpoons H^++CO_3^{2-}}
\]
である。CO2の増加でpHが低下する酸性化は、海水が直ちに酸性のpH7未満になることを意味しない。溶存無機炭素、pH、全アルカリ度は異なる変数である。海洋への人為的炭素の蓄積は大気の増加を和らげる一方、海洋化学の状態を変える。全球の正味吸収と局所的な放出域を区別する。\cite{noaa_carbon_chemistry}

\subsection{潮汐は場所ごとの応答を含む}
\label{theory:ocean:tides}
月と太陽の引力の空間差が海洋を周期的に駆動する。月の半日周の主要成分は約12時間25分、太陽の半日周成分は約12時間である。1日に満潮と干潮が二つずつという半日周潮に加え、日周潮と混合潮もある。原典の「通常2回」は代表像であり、すべての海岸の厳密な頻度ではない。地形、海底、水の慣性と回転による局所応答を含めて潮位を測る。\cite{noaa_tidal_patterns}

\subsection{長波として津波速度を導く}
\label{theory:ocean:tsunami}
水深\(h\)、小さい海面変位\(\eta\)、水深平均水平流速\(u\)を考える。非粘性、一定水深、波長が水深より長く、振幅が水深より小さい一方向の線形浅水モデルは
\[
 \partial_t\eta+h\partial_xu=0,\qquad
 \partial_tu+g\partial_x\eta=0.
\]
二式から\(\partial_t^2\eta-gh\partial_x^2\eta=0\)となり、波速は\(c=\sqrt{gh}\)である。水深の二乗に比例しない。例えば\(h=4000\)mなら\(c\simeq198\)m/s、約713 km/hである。これは一様水深のモデル計算であり実経路の予報値ではない。\cite{jma_tsunami_prediction}

\subsection{浅水変形・湾の集束・河川遡上}
\label{theory:ocean:shoaling}
反射・砕波・散逸が小さく、ゆっくり変化する水深、一定波面幅を仮定する。振幅\(a\)の長波のエネルギー密度は\(a^2\)に比例するため、流束保存\(a^2c\simeq\text{一定}\)と\(c\propto h^{1/2}\)から\(a\propto h^{-1/4}\)を得る。これがグリーンの法則の単純な形である。湾の幅、反射と共振、海岸の砕波があれば別の効果が重なる。\cite{jma_tsunami_prediction}
V字湾の奥では波が集まって増幅し得る。川へ入った津波も陸の内部まで進むことがある。海岸からの距離だけで判定せず、河道と地形を考える。\cite{jma_tsunami_generation,jma_tsunami_rivers}

\subsection{過去問に戻るときの確認項目}
塩分は単位と分母、環流は半球と種類、栄養塩は取込みと再生、堆積物は供給と溶解、津波は速度と高さを先に分ける。原典の「最大」「どこも」「深海ほど」を一般傾向から全称命題へ変えない。頻出度と学習優先度の最終集計は全問の確認後に行う。
